\section{La autoevolución del CEO: valentía y sabiduría}

\subsection{Visión del talento: Hungry, Humble, Smart, Clarity}

Las cuatro características que Hu Yuanming (胡渊鸣) más valora al entrevistar a cada candidato de Meshy:

\begin{importantbox}{El modelo de cuatro dimensiones del talento de Hu Yuanming}
\begin{enumerate}
\item \textbf{Hungry}: aún no se ha demostrado a sí mismo, tiene muchas ganas de convertirse en una mejor versión de sí mismo, siente una especie de "hambre" por generar impact. Quien ya alcanzó el éxito y el reconocimiento puede no estar dispuesto a take risk.
\item \textbf{Humble}: es capaz de reconocer que todavía tiene muchas cosas que no ha hecho bien, tiene growth mindset, y se transforma y se cuestiona a sí mismo constantemente. "Si una persona no es humble, por muy bueno que sea en este momento, en el futuro no avanzará."
\item \textbf{Smart}: el mejor problem solver, aquel que, con recursos limitados, encuentra todas las formas posibles de resolver un problema (por ejemplo, entrenar con números de punto flotante de menor precisión, o encontrar un mejor optimizer).
\item \textbf{Clarity}: expresarse con claridad. "Cuando te das cuenta de que no entiendes en absoluto lo que dice una persona, en nueve de cada diez casos el problema no eres tú, sino que esa persona no ha pensado las cosas con claridad." Expresarse de forma poco clara casi siempre no es un problema de elocuencia, sino de pensamiento poco claro.
\end{enumerate}
\end{importantbox}

Él admite que también comete errores frecuentemente en estos cuatro aspectos. Uno de los cambios más grandes tiene que ver con Clarity: pasar de "decir mucho pero que solo el 10\% sea útil" a darse cuenta de que "escuchar es mucho más importante que hablar". En cuanto a Smart, el error más grande fue "be smart about the solution, but ignore being smart about strategy": pasar demasiado tiempo pensando en el how y muy poco tiempo pensando en el why.

\subsection{La brecha cognitiva entre ser Intern y ser CEO}

Toda la experiencia profesional de Hu Yuanming antes de ser CEO consistió en haber sido becario en NVIDIA, Adobe y Microsoft. "Mi experiencia más professional fue haber sido intern: eso era todo mi conocimiento del mundo laboral antes de ser CEO."

El struggle inicial no era "cómo ser un buen CEO", sino "ni siquiera sé cómo debería ser una empresa normal". Además, como el financiamiento de Taichi (太极) fue muy fluido, no existía la presión de sobrevivencia de "el próximo mes no podré pagar los salarios", y esto en realidad delay el momento en que tuvo que enfrentar los hechos crueles: "que el financiamiento fluya bien en realidad es un arma de doble filo; te hace sentir menos dolor, y entonces no sientes suficiente urgent para cambiar".

\begin{knowledgebox}{Un PhD no es condición suficiente para ser CEO}
"Un PhD te puede enseñar a ser humble, work hard, presentation skills, tecnología de vanguardia. Pero el mercado, los usuarios, la producción en volumen de un producto: eso el PhD no lo toca." Descubrió que ejecutivos de grandes empresas veinte años mayores que él, que salieron a emprender, "parece que lo hicieron incluso peor que nosotros"; esto demuestra que no es un problema de experiencia, sino de velocidad de aprendizaje y de qué tan humble se es. Incluso Li Xiang (李想), de Li Auto (理想汽车), que ni siquiera terminó la preparatoria, lo hizo muy bien: "un PhD no es condición necesaria, la visión sí lo es."
\end{knowledgebox}

Se autoevalúa como "so far un CEO que apenas aprueba, a duras penas", y entre los aspectos que necesita mejorar están: delegate mejor el trabajo, administrar el tiempo, ser más decidido en ciertas cosas, aprender más rápido, dedicar más tiempo a lo externo que a lo interno, leer más libros, reflexionar durante más tiempo, y pensar en la visión de la empresa.

\subsection{La cultura de ganar y la cultura de perder}

\begin{warningbox}{Características de la cultura de perder}
\begin{itemize}
\item \textbf{Desgaste interno}: los departamentos se desgastan entre sí en lugar de cooperar
\item \textbf{Pasar la responsabilidad a otros}: cuando surge un problema, no se revisa lo ocurrido ni se reflexiona
\item \textbf{Poco realista}: no se trata de que la meta sea alta (la meta puede seguir siendo alta), sino de no enfrentar los hechos: disfrazar lo malo como si fuera bueno
\item \textbf{Buscar la armonía}: el equipo no puede criticar al CEO, no puede señalar los problemas en público
\end{itemize}
Todo esto termina llevando a la empresa a "perder".
\end{warningbox}

Por el contrario, el pilar cultural de Meshy es "Respect Facts" (respetar los hechos). Hu Yuanming dio un ejemplo concreto: hace unos días, dos colegas lo "llevaron al pasillo" para criticarle problemas de gestión. "Cuando lo escuché se me encogió el corazón, pero dejé que terminara de hablar, y luego pensé si de verdad no lo había hecho bien; después descubrí que efectivamente había cosas que no había hecho bien."

Él considera que poder contratar a este tipo de personas y tener este tipo de ambiente ya es en sí una suerte: "Esa persona seguramente tiene un fuerte sentido de ownership hacia la empresa, se ve a sí misma casi como el CEO, y por eso está dispuesta a decirte estas cosas."

En una empresa impulsada por la innovación, el CEO no necesariamente sabe todo: esta es la diferencia esencial con una empresa manufacturera. "En una empresa impulsada por la innovación tienes que aceptar que el CEO no necesariamente sabe todo."

\subsection{La experiencia de trabajar en Meshy}

Empleados veteranos con tres o cuatro años de antigüedad describen a Meshy como "la mejor empresa en la que he trabajado en más de diez años de carrera". Las razones incluyen: poder crecer rápidamente, tener un leader que escucha sus necesidades, desempeñar un papel importante en un negocio de rápido crecimiento, y tener acceso a la tecnología más de vanguardia (antes, en una gran empresa, tal vez solo escribían shaders, mientras que en Meshy pueden entrenar modelos con 512 GPU).

"En este proceso, todos se han convertido en una mejor versión de sí mismos." Esto coincide en gran medida con su objetivo como CEO: hacer que Meshy sea "un lugar que la gente más talentosa considere un excelente lugar para trabajar".

\subsection{El cambio de perspectiva más profundo: dejar de buscar que todos lo quieran}

Hu Yuanming dice que este fue su mayor momento de colapso cognitivo como CEO:

\begin{warningbox}{Buscar que todos te quieran es un acto egoísta}
"Si el objetivo de tu trabajo es que todos te quieran a ti, al final terminarás siendo odiado por todos." Porque el CEO tiene que tomar tough decisions: cancelar proyectos que le gustan a todos, despedir a colegas sin contribución pero con buenas relaciones, elegir entre el grupo de usuarios A y el grupo B (si haces lo que le gusta a A, B te critica). "Buscar que todos te quieran es poner lo personal antes que lo colectivo: cuando ocupas la posición de CEO, tienes que responder por los intereses de toda la empresa, no por que todos te quieran a ti."
\end{warningbox}

Citó la evaluación que Li Xiang, fundador de Li Auto (理想汽车), hizo sobre Elon Musk: "Evaluar a un CEO no es ver qué tan diplomático es, sino ver si en un momento de verdadera crisis es capaz de ver la esencia del problema de un vistazo."

Este cambio también se extiende al nivel de gestión. Descubrió que los nuevos managers suelen tener el mismo problema: como fueron tratados con consideración cuando ellos mismos eran IC (individual contributor), tratan a sus subordinados de la misma manera y no se atreven a exigir estándares más altos. La solución es contarles la "historia de sangre y lágrimas": "El problema que tiene hoy la empresa es porque hace medio año fui indeciso y quise que todos me quisieran."

\subsection{Las reglas de supervivencia de una startup}

\textbf{¿Qué es un buen negocio?} Hu Yuanming resumió tres dimensiones:
\begin{enumerate}
\item \textbf{Mercado grande}: la cuestión de si vale la pena hacerlo
\item \textbf{No consenso}: la cuestión de si todavía queda una oportunidad para ti (el timing es muy importante)
\item \textbf{Ventaja competitiva}: la cuestión de si tú mismo puedes lograrlo
\end{enumerate}

\begin{importantbox}{El volante tecnológico (Flywheel)}
Buena tecnología $\rightarrow$ buen producto $\rightarrow$ buen revenue $\rightarrow$ atrae a buena gente $\rightarrow$ produce buena tecnología. Si la tecnología no se puede comercializar, se convierte en "un espectáculo de fuegos artificiales que desaparece en cuanto termina". El volante de Meshy ya está girando: los ingresos mensuales crecen 20\%, cuenta con 4 millones de usuarios, y la tecnología se sigue iterando constantemente.
\end{importantbox}

\textbf{Sobre elegir una dirección}, citó a Andrej Karpathy: "Algo que rinde diez veces más suele ser solo el doble de difícil." Elegir la dirección correcta es mucho más importante que el nivel de esfuerzo: el equipo de Cursor es tan hardworking como el de Meshy, pero el sector en el que compite Cursor tiene un retorno mayor. Sin embargo, él también lo ve de manera dialéctica: "Lo que se obtiene por suerte, tarde o temprano hay que devolverlo con capacidad real. Que Cursor crezca rápido ahora es algo temporal; que en la siguiente etapa pueda crecer aún más dependerá de si el CEO es capaz de seguir aceptando nuevos retos."

\subsection{Respect Facts: el pilar cultural de respetar los hechos}

\begin{warningbox}{El hecho más cruel al emprender}
"La gran mayoría del esfuerzo en el mundo no tiene ninguna recompensa. Mientras más grande es lo que haces y mayor es el riesgo, muchas veces pones 99 puntos de esfuerzo y no obtienes ni un punto de recompensa. Pero tienes que hacer ese esfuerzo, porque no sabes cuál de esos 100 puntos es el que sirve. 'El cielo recompensa a los diligentes', 'un esfuerzo, una cosecha': todo esto es pura mentira."
\end{warningbox}

En la práctica de evaluación de Meshy, exigen pruebas ciegas: muestran de forma aleatoria la versión nueva y la antigua, sin decir cuál es cuál, y un tercero las evalúa. "Muchas veces, después de mucho esfuerzo, descubres que la versión anterior era mejor." Si no se enfrenta este hecho de manera objetiva, el equipo cae en el autoengaño.

\subsection{Bias for Action y empezar con el fin en mente}

Hu Yuanming se describe a sí mismo como alguien "muy bias for action": si quiere hacer algo, lo hace de inmediato, incluso sin dormir por la noche. Con frecuencia, cuando un proyecto todavía es apenas una idea sin forma, ya está pensando en el plano del futuro: "es un poco como Elon Musk construyendo cohetes, que antes de que el motor esté listo ya está pensando qué va a comer en Marte."

Pero esto no es un activismo ciego: él piensa "empezando con el fin en mente", es decir, qué debe hacer ahora para lograr el objetivo final, y luego planea hacia atrás a partir de ahí. Tanto Taichi como Meshy se planearon de esta manera, hacia atrás.

\begin{knowledgebox}{¿Si Jensen Huang volviera a empezar, no fundaría NVIDIA?}
"Creo que eso es alardear fingiendo modestia. Tienes una empresa de 4 billones de dólares; puedes decir lo que quieras, todo lo que digas será correcto, tú eres quien manda." Hu Yuanming cambió de tono: "\textbf{Yo sí volvería a elegirlo.} Aunque cada día es muy suffer, en este proceso se puede aprender mucho y convertirse en una mejor versión de uno mismo. Se ve una visión mucho más amplia que la del pequeño comedor del mundo académico."
\end{knowledgebox}

"En tiempos de paz, no hay nada que forje más a una persona que emprender." Esta es su evaluación final sobre el emprendimiento. No se puede esperar aprender demasiado en un entorno protegido (la escuela, una gran empresa), a menos que tengas un mentor excepcionalmente bueno. Emprender es la forma de crecimiento más cruel, pero también la más eficiente.

\subsection{Resumen de la sección}

La evolución personal del CEO es un proceso continuo de "matar a la versión pasada de uno mismo". Los cambios centrales incluyen: pasar de buscar que todos lo quieran a buscar el éxito del negocio; pasar de smart about solution a smart about strategy; pasar de trabajar con la cabeza agachada a abrir la mirada y comunicarse más. Las características de una cultura que puede ganar son enfrentar la realidad, aceptar la crítica y buscar la excelencia; las características de una cultura que puede perder son el desgaste interno, pasar la responsabilidad a otros y disfrazar los problemas.
